% Part: propositional-logic
% Chapter: propositional-logic 
% Section: completeness

\documentclass[../../../include/open-logic-section]{subfiles}

\begin{document}

\olfileid[pt]{pl}{prp}{com}

\olsection{Completude da lógica proposicional}

\begin{defn}
\ollabel{def:MaxCon}
Um conjunto $\Gamma$ de !!{formula}s é \emph{maximalmente consistente}
se é consistente e, sempre que $\Delta$ é um conjunto consistente tal
que $\Gamma \subseteq \Delta$, temos $\Gamma = \Delta$.
\end{defn}

\begin{lem}[Lema da verdade]
  \ollabel{lem:truth}
Seja $\Gamma$ maximalmente consistente; então:
\begin{enumerate}
\item \ollabel{lem:truth-1} $!A \in \Gamma$ se, e somente se, $\lnot !A
  \notin \Gamma$;
\item \ollabel{lem:truth-2} $\Gamma \Proves !A$ se, e somente se, $!A \in \Gamma$; 
\item \ollabel{lem:truth-3} $!A \lif !B \in \Gamma$ se, e somente se,
  $!A \notin \Gamma$ ou $!B \in \Gamma$.
\end{enumerate}
\end{lem}

\begin{proof}
\begin{enumerate}
\item Suponha $!A \in \Gamma$. Se também $\lnot !A \in \Gamma$, então
  $\Gamma$ é inconsistente; e, se nem $!A$ nem $\lnot!A$ pertencem a
  $\Gamma$, então, por \olref[axd]{prop:phi}, um dos conjuntos
  $\Gamma\cup\{!A\}$ e $\Gamma \cup \{\lnot !A\}$ é consistente, o que
  significa que $\Gamma$ não é maximal. A outra direção é demonstrada
  da mesma maneira.
\item Se $!A \in \Gamma$, então, naturalmente, $\Gamma \Proves !A$.
  Por outro lado, suponha $\Gamma \Proves !A$. Se $!A \notin \Gamma$,
  então, por \olref{lem:truth-1}, $\lnot !A \in \Gamma$. Logo,
  $\Gamma \Proves !A$ e $\Gamma \Proves \lnot !A$, uma contradição.
\item Exercício.
\end{enumerate}
\end{proof}

A direção da esquerda para a direita de \olref{lem:truth}\olref{lem:truth-2}
mostra que $\Gamma$ é \emph{dedutivamente fechado}, O.s.\@, se
$\Gamma \Proves !A$, então $!A \in \Gamma$ para todo~$!A$.

\begin{prob}
  Complete a demonstração de \olref{lem:truth}.
\end{prob}

\begin{thm}[Completude] 
\ollabel{thm:completeness} 
Se $\Gamma$ é consistente, então $\Gamma$ é satisfatível.
\end{thm}

\begin{proof}
Seja $!A_0$, $!A_1$,~\dots uma enumeração exaustiva de todas as
!!{formula}s da linguagem. Definimos recursivamente uma sequência
crescente de conjuntos de !!{formula}s $\Gamma_0$, $\Gamma_1$,~\dots,
estabelecendo:
\begin{align*}
\Gamma_0 & = \Gamma\\
\Gamma_{n+1} &  =
\begin{cases}
  \Gamma_n \cup \{!A_n\} & \text{se $\Gamma_n \cup \{!A_n\}$ é
    consistente;}\\
  \Gamma_n \cup \{\lnot !A_n\} & \text{caso contrário}.
\end{cases}
\end{align*}
Defina, então:
\[
\Gamma^* = \bigcup_{0\le n}\Gamma_n.
\]
A demonstração prossegue estabelecendo, na ordem, os seguintes fatos:
\begin{enumerate}
\item Para cada $n$, o conjunto $\Gamma_n$ é consistente (por indução
  sobre $n$, usando \olref[axd]{prop:phi});
\item $\Gamma^*$ é consistente;
\item $\Gamma^*$ é maximal.
\end{enumerate}
Defina uma valoração $v$ estabelecendo $v(p_i) = \True$ se, e somente
se, $p_i \in \Gamma^*$. Por indução sobre $!A$, mostra-se então que a
pertença a $\Gamma^*$ coincide com a verdade segundo $v$ também para
!!{sentence}s mais complexas: $\overline{v}(!A) = \True$ se, e somente
se, $!A \in \Gamma^*$. Em particular, $v$ satisfaz $\Gamma^*$ e, como
$\Gamma \subseteq \Gamma^*$, $\Gamma$ também é satisfatível, como
queríamos.
\end{proof}

\begin{cor}
Se $\Gamma \Entails !A$, então $\Gamma \Proves !A$.
\end{cor}

\begin{proof}
Se $\Gamma\Proves/ !A$, então $\Gamma \cup \{\lnot !A\}$ é consistente,
por \olref[axd]{prop:prov-incons}. Portanto, pelo teorema da completude,
$\Gamma \cup \{\lnot !A\}$ é satisfatível, e $\Gamma \Entails/ !A$.
\end{proof}

\begin{prop}[Teorema da compacidade]
\ollabel{prop:compactness}
$\Gamma$ é satisfatível se, e somente se, todo subconjunto \emph{finito}
$\Gamma_0$ de $\Gamma$ é satisfatível.
\end{prop}

\begin{proof}
$\Gamma$ é insatisfatível se, e somente se, é inconsistente; se, e
  somente se, algum subconjunto finito $\Gamma_0$ de $\Gamma$ é
  inconsistente; se, e somente se, algum subconjunto finito $\Gamma_0$
  de $\Gamma$ é insatisfatível.
\end{proof}

\begin{cor}
$\Gamma \Entails !A$ se, e somente se, para algum subconjunto finito
$\Gamma_0$ de $\Gamma$, $\Gamma_0 \Entails !A$.
\end{cor}

\end{document}
